% Versión LaTeX de docs/metodologia.md (2026-11-14). La fuente es el .md: si cambia, actualizar aquí.
\documentclass[11pt,a4paper]{article}
% Preámbulo común de la documentación LaTeX de TerraPulse AI.
% Se incluye con % Preámbulo común de la documentación LaTeX de TerraPulse AI.
% Se incluye con % Preámbulo común de la documentación LaTeX de TerraPulse AI.
% Se incluye con \input{preambulo} desde cada documento. Compilar con pdflatex (o latexmk -pdf).
\usepackage[utf8]{inputenc}
\usepackage[T1]{fontenc}
\usepackage{lmodern}
\usepackage[spanish,es-noshorthands,es-tabla,es-nodecimaldot]{babel}
\usepackage{amsmath,amssymb}
\usepackage[margin=2.3cm]{geometry}
\usepackage{array,booktabs,longtable}
\usepackage[normalem]{ulem}
\usepackage{xcolor}
\usepackage{enumitem}
\usepackage{listings}
\usepackage[hidelinks]{hyperref}

\setlength{\parindent}{0pt}
\setlength{\parskip}{0.5em}
\setlist{itemsep=2pt,topsep=2pt}
\renewcommand{\arraystretch}{1.25}
\Urlmuskip=0mu plus 1mu

% Símbolos frecuentes
\newcommand{\sz}{\ensuremath{\sigma^0}}
\newcommand{\gz}{\ensuremath{\gamma^0}}
\newcommand{\bz}{\ensuremath{\beta^0}}
\newcommand{\dsz}{\ensuremath{\Delta\sigma^0}}
\newcommand{\ok}{\ensuremath{\checkmark}}
\newcommand{\flecha}{\ensuremath{\rightarrow}}
\newcommand{\aprox}{\ensuremath{\approx}}
\newcommand{\por}{\ensuremath{\times}}

% Recuadro de nota (equivale a los ">" del Markdown)
\newenvironment{nota}
  {\begin{quote}\small\itshape}
  {\end{quote}}

% Bloques de código y comandos
\lstset{
  basicstyle=\ttfamily\small,
  breaklines=true,
  columns=fullflexible,
  keepspaces=true,
  frame=single,
  rulecolor=\color{black!30},
  backgroundcolor=\color{black!3},
  xleftmargin=4pt, xrightmargin=4pt,
  extendedchars=true,
  literate={á}{{\'a}}1 {é}{{\'e}}1 {í}{{\'i}}1 {ó}{{\'o}}1 {ú}{{\'u}}1
           {Á}{{\'A}}1 {É}{{\'E}}1 {Í}{{\'I}}1 {Ó}{{\'O}}1 {Ú}{{\'U}}1
           {ñ}{{\~n}}1 {Ñ}{{\~N}}1
}
 desde cada documento. Compilar con pdflatex (o latexmk -pdf).
\usepackage[utf8]{inputenc}
\usepackage[T1]{fontenc}
\usepackage{lmodern}
\usepackage[spanish,es-noshorthands,es-tabla,es-nodecimaldot]{babel}
\usepackage{amsmath,amssymb}
\usepackage[margin=2.3cm]{geometry}
\usepackage{array,booktabs,longtable}
\usepackage[normalem]{ulem}
\usepackage{xcolor}
\usepackage{enumitem}
\usepackage{listings}
\usepackage[hidelinks]{hyperref}

\setlength{\parindent}{0pt}
\setlength{\parskip}{0.5em}
\setlist{itemsep=2pt,topsep=2pt}
\renewcommand{\arraystretch}{1.25}
\Urlmuskip=0mu plus 1mu

% Símbolos frecuentes
\newcommand{\sz}{\ensuremath{\sigma^0}}
\newcommand{\gz}{\ensuremath{\gamma^0}}
\newcommand{\bz}{\ensuremath{\beta^0}}
\newcommand{\dsz}{\ensuremath{\Delta\sigma^0}}
\newcommand{\ok}{\ensuremath{\checkmark}}
\newcommand{\flecha}{\ensuremath{\rightarrow}}
\newcommand{\aprox}{\ensuremath{\approx}}
\newcommand{\por}{\ensuremath{\times}}

% Recuadro de nota (equivale a los ">" del Markdown)
\newenvironment{nota}
  {\begin{quote}\small\itshape}
  {\end{quote}}

% Bloques de código y comandos
\lstset{
  basicstyle=\ttfamily\small,
  breaklines=true,
  columns=fullflexible,
  keepspaces=true,
  frame=single,
  rulecolor=\color{black!30},
  backgroundcolor=\color{black!3},
  xleftmargin=4pt, xrightmargin=4pt,
  extendedchars=true,
  literate={á}{{\'a}}1 {é}{{\'e}}1 {í}{{\'i}}1 {ó}{{\'o}}1 {ú}{{\'u}}1
           {Á}{{\'A}}1 {É}{{\'E}}1 {Í}{{\'I}}1 {Ó}{{\'O}}1 {Ú}{{\'U}}1
           {ñ}{{\~n}}1 {Ñ}{{\~N}}1
}
 desde cada documento. Compilar con pdflatex (o latexmk -pdf).
\usepackage[utf8]{inputenc}
\usepackage[T1]{fontenc}
\usepackage{lmodern}
\usepackage[spanish,es-noshorthands,es-tabla,es-nodecimaldot]{babel}
\usepackage{amsmath,amssymb}
\usepackage[margin=2.3cm]{geometry}
\usepackage{array,booktabs,longtable}
\usepackage[normalem]{ulem}
\usepackage{xcolor}
\usepackage{enumitem}
\usepackage{listings}
\usepackage[hidelinks]{hyperref}

\setlength{\parindent}{0pt}
\setlength{\parskip}{0.5em}
\setlist{itemsep=2pt,topsep=2pt}
\renewcommand{\arraystretch}{1.25}
\Urlmuskip=0mu plus 1mu

% Símbolos frecuentes
\newcommand{\sz}{\ensuremath{\sigma^0}}
\newcommand{\gz}{\ensuremath{\gamma^0}}
\newcommand{\bz}{\ensuremath{\beta^0}}
\newcommand{\dsz}{\ensuremath{\Delta\sigma^0}}
\newcommand{\ok}{\ensuremath{\checkmark}}
\newcommand{\flecha}{\ensuremath{\rightarrow}}
\newcommand{\aprox}{\ensuremath{\approx}}
\newcommand{\por}{\ensuremath{\times}}

% Recuadro de nota (equivale a los ">" del Markdown)
\newenvironment{nota}
  {\begin{quote}\small\itshape}
  {\end{quote}}

% Bloques de código y comandos
\lstset{
  basicstyle=\ttfamily\small,
  breaklines=true,
  columns=fullflexible,
  keepspaces=true,
  frame=single,
  rulecolor=\color{black!30},
  backgroundcolor=\color{black!3},
  xleftmargin=4pt, xrightmargin=4pt,
  extendedchars=true,
  literate={á}{{\'a}}1 {é}{{\'e}}1 {í}{{\'i}}1 {ó}{{\'o}}1 {ú}{{\'u}}1
           {Á}{{\'A}}1 {É}{{\'E}}1 {Í}{{\'I}}1 {Ó}{{\'O}}1 {Ú}{{\'U}}1
           {ñ}{{\~n}}1 {Ñ}{{\~N}}1
}


\title{Metodología}
\author{TerraPulse AI}
\date{Borrador, 14 de noviembre de 2026}

\begin{document}
\maketitle

\begin{nota}
\textbf{Borrador.} Se documenta conforme se ejecuta el pipeline (\texttt{AGENTS.md} §6). Dueño: P3.
Todo lo que aquí aparece como «verificado» se comprobó con el comando indicado y la fecha anotada;
lo que aún no se ha comprobado se marca como \textbf{pendiente}.
\end{nota}

\section{Región de interés (ROI)}

Chiapas, México. Polígono en EPSG:4326 (\path{pipeline/config.yaml -> roi}):

\begin{lstlisting}
POLYGON((-93.3919 16.6539,-92.8757 16.6539,-92.8757 16.8964,-93.3919 16.8964,-93.3919 16.6539))
\end{lstlisting}

\section{Escenas utilizadas}

Fuente: catálogo de ASF consultado con \path{python -m pipeline.catalog} el 2026-11-14
(\path{data/catalog/catalog_asf_GCOV_20261114T231456Z.json}). Recortes obtenidos con
\path{python -m pipeline.download} (lectura remota por ventanas; solo se transfieren los bloques de la ROI).

\textbf{Par de trabajo (Fase 1):}

\begin{longtable}{p{0.30\linewidth}p{0.31\linewidth}p{0.31\linewidth}}
\toprule
 & \textbf{t1} & \textbf{t2} \\
\midrule
\endhead
Producto & NISAR L2 GCOV & NISAR L2 GCOV \\
Fecha & 2026-06-18 & 2026-10-16 \\
Trayectoria / marco / dirección & 62 / 10 / ASCENDING & 62 / 10 / ASCENDING \\
Polarizaciones (Frequency A, leídas del .h5) & HH, HV & HH, HV \\
CRID / madurez & P05023 / PROVISIONAL & P05023 / PROVISIONAL \\
Píxeles válidos en el polígono ROI (capa \texttt{mask}) & 99.57\,\% & 99.61\,\% \\
Escena & \path{NISAR_L2_PR_GCOV_023_062_A_010_4005_DHDH_A_20260618T114010_20260618T114044_P05023_N_F_J_001}
       & \path{NISAR_L2_PR_GCOV_033_062_A_010_4005_DHDH_A_20261016T114006_20261016T114041_P05023_N_F_J_001} \\
\bottomrule
\end{longtable}

Cada recorte está en \path{data/raw/<scene_id>/} con su \path{scene.json} (procedencia, URL, ventana, malla, cobertura).

\textbf{Por qué T62 y no T55.} La pila T55/F80 DESCENDING P05023 tiene las mismas fechas, pero su marco
solo cubre la mitad este de la ROI: 50.3\,\% de píxeles válidos (32.6\,\% fuera de la adquisición,
17.1\,\% no válidos), con datos solo al este de lon \aprox{} $-93.14$. Se conserva en \path{data/raw/} como respaldo.
La pila T156/F80 no se ha recortado (cobertura \textbf{pendiente}).

\textbf{Escena sin pareja.} GCOV T62/F10 2026-10-28 (P05036, VALIDATED) se inspeccionó en la Etapa 1
(\path{data/explore/}), pero es la única fecha de su pila: el CRID cambió de P05023 a P05036 el 2026-10-23
y no se mezclan versiones de procesamiento.

\section{Comparabilidad}

Verificada para el par de trabajo con \path{compat.check_pair} (catálogo) y \path{sources/subset.compare_grids} (.h5):

\begin{itemize}
  \item Mismo producto, trayectoria, marco, dirección, polarizaciones, CRID y madurez.
  \item Misma malla en ambas fechas: EPSG:32615 (UTM 15N), 10\,\por\,10\,m, recorte de 2687\,\por\,5504\,px con
        idéntica transformación geográfica. No se requiere remuestreo.
  \item La identidad de cada archivo (granuleId, CRID, trayectoria, marco, dirección) se contrastó con el catálogo al recortar.
\end{itemize}

\section{Preprocesamiento y máscaras}

\textbf{Pendiente} (\path{preprocess.py}, \path{masks.py}, P2). Datos de partida y advertencias que declara el producto:

\begin{itemize}
  \item Las capas \texttt{HHHH} y \texttt{HVHV} son potencia lineal \gz{} (convención leída del .h5).
        $\sigma^0 = \gamma^0 \cdot \texttt{rtcGammaToSigmaFactor}$.
  \item \path{noiseCorrectionApplied = False}: valores bajos, sobre todo en HV, pueden estar cerca del piso de ruido.
  \item \path{shadowMaskingApplied = False}: la máscara de sombra/layover se debe aplicar aparte.
  \item Capas de potencia guardadas con mantisa de 16 bits (\textasciitilde 5 cifras significativas).
  \item Los recortes se entregan sin dB, sin máscara y sin multilook; es el rectángulo que contiene la ROI.
\end{itemize}

\section{Métrica de cambio}

\textbf{Pendiente} (\path{change.py}, P2). Prevista:
\[
  \Delta\sigma^0_{\text{dB}} = \sigma^0_{\text{dB},t_2} - \sigma^0_{\text{dB},t_1},
\]
razón y z-score respecto a la pila.
Un cambio en la retrodispersión describe un cambio en la respuesta radar; no identifica por sí solo su causa.
Las ecuaciones completas están en \texttt{ecuaciones\_radar.tex}.

\section{Validación}

\textbf{Pendiente: depende del fenómeno, que aún no se decide.} Fuentes auxiliares candidatas encontradas
en NASA Earthdata (CMR) para la ROI entre 2026-06-01 y 2026-10-31. Consulta del 2026-11-14,
\textbf{solo búsqueda en el catálogo}: aún no se ha descargado ni inspeccionado su contenido sobre la ROI.

{\small
\begin{longtable}{p{0.20\linewidth}p{0.21\linewidth}p{0.12\linewidth}p{0.10\linewidth}p{0.26\linewidth}}
\toprule
\textbf{Producto (short name)} & \textbf{Qué es} & \textbf{Resolución} & \textbf{Archivo (DAAC)} & \textbf{Disponibilidad sobre la ROI} \\
\midrule
\endhead
\path{GPM_3IMERGDL} v07 & Precipitación diaria GPM IMERG (Late) & 0.1° (\textasciitilde 11\,km), 1 día & GES DISC & 153 días (01-jun \flecha{} 31-oct) \\
\path{GPM_3IMERGDF} v07 & Precipitación diaria GPM IMERG (Final) & 0.1°, 1 día & GES DISC & 0: aún no publicada para 2026 \\
\path{HLSS30} v2.0 & Reflectancia de superficie Sentinel-2 armonizada & 30\,m & LP DAAC & 64 fechas; 9 con nubes $\leq$ 20\,\% en el tile: 06-20, 06-30, 07-20, 08-01, 08-14, 08-17, 08-29, 09-30, 10-01 \\
\path{HLSL30} v2.0 & Reflectancia de superficie Landsat armonizada & 30\,m & LP DAAC & 18 fechas; 6 con nubes $\leq$ 20\,\% en el tile: 06-21, 07-23, 08-08, 08-16, 09-01, 10-27 \\
\path{OPERA_L3_DSWX-HLS_V1} & Extensión de agua superficial (óptico) & 30\,m & PO.DAAC & 80 fechas \\
\path{OPERA_L3_DSWX-S1_V1} & Extensión de agua superficial (Sentinel-1, banda C) & 30\,m & PO.DAAC & 65 fechas \\
\path{OPERA_L3_DIST-ALERT-HLS_V1} & Alertas de perturbación de la vegetación (óptico) & 30\,m & LP DAAC & 74 fechas \\
\path{OPERA_L3_DIST-ALERT-S1_V1} & Alertas de perturbación (Sentinel-1) & 30\,m & --- & 0 \\
\bottomrule
\end{longtable}
}

Notas:

\begin{itemize}
  \item El porcentaje de nubes de HLS es del tile completo, no de la ROI: comprobarlo sobre la ROI antes de usar una fecha.
  \item Fechas ópticas cercanas al par: 2026-06-20 / 06-21 (t1 = 06-18) y 2026-10-01 / 10-27 (t2 = 10-16).
  \item IMERG (\textasciitilde 11\,km) sirve como \textbf{contexto} temporal (lluvia entre fechas y antes de cada adquisición),
        no como validación píxel a píxel.
\end{itemize}

\textbf{Fuente auxiliar sugerida según el fenómeno} (decisión de P3):

\begin{longtable}{p{0.25\linewidth}p{0.45\linewidth}p{0.20\linewidth}}
\toprule
\textbf{Fenómeno} & \textbf{Referencia} & \textbf{Contexto} \\
\midrule
\endhead
Inundación / humedales & OPERA DSWx (HLS y/o S1) & GPM IMERG \\
Vegetación / perturbación & OPERA DIST-ALERT-HLS, HLS (índices de vegetación) & GPM IMERG \\
Deformación & Ninguna de las anteriores & --- \\
\bottomrule
\end{longtable}

\textbf{Precaución.} OPERA es otro producto derivado de satélite, con su propio error; no son etiquetas de campo.
Precisión, recall, F1 e IoU calculados contra OPERA miden el \textbf{acuerdo con OPERA}, no la exactitud real,
y así deben reportarse. Sin una referencia adecuada, el resultado se marca como \textbf{exploratorio}.

\subsection{Sensores evaluados y descartados (por ahora)}

Recursos de agencias socias incluidos en la lista del reto:

\begin{longtable}{p{0.17\linewidth}p{0.07\linewidth}p{0.20\linewidth}p{0.46\linewidth}}
\toprule
\textbf{Sensor} & \textbf{Banda} & \textbf{Decisión} & \textbf{Motivo} \\
\midrule
\endhead
RADARSAT-1 (CSA) & C & Descartado & Operó de 1995 a 2013: sin datos para 2026 \\
RCM (CSA) & C & Descartado por ahora & Otra banda y otra calibración: no comparable con NISAR; condiciones de acceso por verificar \\
SAOCOM (CONAE) & L & Posible contraste en la Fase 2 & Misma banda, pero otro sensor: nunca se mezcla en un \dsz{} con NISAR; cobertura de la ROI y fechas \textbf{pendientes} de consultar (requiere registro) \\
\bottomrule
\end{longtable}

Un segundo radar no confirma la causa física de un cambio de retrodispersión; la validación requiere datos
de otra naturaleza (\texttt{AGENTS.md} §6).

\section{Limitaciones}

\begin{itemize}
  \item Productos L-band PROVISIONAL: calibrados y parcialmente validados, no completamente validados.
  \item Sin corrección de ruido ni máscara de sombra/layover en el producto (ver Preprocesamiento).
  \item Un solo par y una sola geometría (ascendente); la pila tiene 10 fechas disponibles para una línea base.
  \item Fenómeno aún no decidido: hasta entonces, cualquier mapa de cambio es exploratorio.
\end{itemize}

\end{document}
